% Reward Hacking Detection via Behavioral Probes
% RH-only extract from paper.tex — no EM, no Neel Q's

\documentclass[11pt]{article}

\usepackage[margin=1in]{geometry}
\usepackage{graphicx}
\usepackage{booktabs}
\usepackage{amsmath}
\usepackage{amssymb}
\usepackage{hyperref}
\usepackage{xcolor}
\usepackage{caption}
\usepackage{subcaption}
\usepackage[numbers]{natbib}

\graphicspath{{figures/}}

\title{Reward Hacking Detection via Behavioral Probes}
\author{Sriram and Eric}
\date{}

\begin{document}
\maketitle

% ============================================================================
\section{Introduction}
% ============================================================================

\begin{itemize}
  \item \textbf{The problem}: can we predict behavior of a fine-tuned model before it's visible in output?
  \begin{itemize}
    \item Outputs aren't enough --- by the time you see misalignment, the decision already happened
    \item CoT can be hidden or faked
    \item Activations are where decisions happen
  \end{itemize}

  \item \textbf{PSM endorsement}
  \begin{itemize}
    \item Anthropic's Persona Selection Model (Marks, Lindsey, Olah --- Feb 2026):
    \begin{itemize}
      \item LLMs simulate personas from pretraining character archetypes
      \item Post-training selects/refines the Assistant persona
      \item Dangerous behaviors use same conceptual vocabulary as pretraining personas
    \end{itemize}
    \item Literal quote: ``as Anthropic does during pre-deployment alignment audits --- build and monitor activation probes for a researcher-curated set of traits like deception and evaluation awareness''
    \item Our work operationalizes this recommendation
  \end{itemize}

  \item \textbf{What we built}: $\sim$150-probe behavioral monitor
  \begin{itemize}
    \item 150 from psychological taxonomy + 20 alignment-specific
    \item Each probe: extracted from base model, causally validated via steering on instruct
    \item One scalar per named behavioral dimension per token
  \end{itemize}

  \item \textbf{Claims}
  \begin{itemize}
    \item \textbf{Claim 1}: Pre-trained behavioral probes detect persona shifts in reward-hacking models
    \item \textbf{Claim 2}: Probes detect persona changes before they're visible in model outputs
    \item \textbf{Claim 3}: Multi-probe decomposition reveals \emph{which} psychological dimensions shift --- a structural fingerprint, not binary detection
  \end{itemize}
\end{itemize}


% ============================================================================
\section{Method}
% ============================================================================

\subsection{Trait Vector Extraction}

\begin{itemize}
  \item \textbf{Taxonomy}: $\sim$150 traits total
  \begin{itemize}
    \item 150 from psychological frameworks (Cowen \& Keltner 2017 27 emotions, Cattell's 16PF, Big Five facets, moral foundations, interpersonal circumplex)
    \item 20 alignment-specific (ulterior\_motive, eval\_awareness, alignment\_faking, \ldots)
    \item Goal: taxonomic coverage --- more probes = more surface area for unsupervised detection
  \end{itemize}

  \item \textbf{Definitions}: per-trait scoring rubric
  \begin{itemize}
    \item What trait IS, what it ISN'T, what's orthogonal
    \item Internal state (felt emotion, conviction) not surface markers
    \item Exhibiting $>$ verbalizing --- expressing anger outweighs saying ``I feel angry''
    \item Reused across pipeline: scenario vetting, steering scoring, fingerprint judge
  \end{itemize}

  \item \textbf{Scenario creation}: contrastive pairs
  \begin{itemize}
    \item First-person document completion prefixes (2.5$\times$ stronger steering than 3rd person)
    \item Hanging pairs at emotional peak (e.g., ``I wanted her to'')
    \item $\sim$15 positive / negative matched pairs per trait
    \item Iterative: reflection + option to recurse at each step
  \end{itemize}

  \item \textbf{Base model completions}
  \begin{itemize}
    \item Base model (no instruct, no chat template) completes each scenario prefix
    \item Why base? More fundamental --- no instruct persona or safety training confounds
    \item PSM: concepts pre-exist alignment. Fine-tuning changes WHEN they activate, not WHERE they live
  \end{itemize}

  \item \textbf{Activation capture}
  \begin{itemize}
    \item Position: \texttt{response[:5]} --- mean of first 5 response tokens
    \begin{itemize}
      \item Base model tends toward neutrality after $\sim$5 tokens
      \item Signal concentrated early (sycophancy peaks at \texttt{[:5]})
    \end{itemize}
    \item Layers: sweep 30--60\% of model depth
    \item Component: residual stream (full layer output)
    \item Output per sample: $\bar{h} = \frac{1}{5}\sum_{t=1}^{5} h_t \in \mathbb{R}^{d}$
  \end{itemize}

  \item \textbf{Vector extraction}: logistic regression (probe)
  \begin{itemize}
    \item Row-normalize each sample to unit norm:
    \begin{equation}
      \tilde{x}_i = \frac{x_i}{\|x_i\|_2}
    \end{equation}
    \item Fit L2-regularized logistic regression ($C=1.0$, \texttt{lbfgs}):
    \begin{equation}
      \hat{w} = \arg\min_w \sum_i \log(1 + e^{-y_i \cdot w^\top \tilde{x}_i}) + \frac{1}{2C}\|w\|_2^2
    \end{equation}
    \item Trait vector = normalized coefficients:
    \begin{equation}
      v = \frac{\hat{w}}{\|\hat{w}\|_2}
    \end{equation}
    \item Why probe over mean\_diff? Probe empirically more robust on short, variable-length responses
  \end{itemize}

  \item \textbf{Steering validation} (causal proof of transfer)
  \begin{itemize}
    \item Apply vector to \emph{instruct} model during generation:
    \begin{equation}
      h_l' = h_l + \alpha \cdot v
    \end{equation}
    \item Layer sweep: 30--60\% depth
    \item Adaptive coefficient search (5 steps):
    \begin{itemize}
      \item Vector is unit normalized ($\|v\| = 1$)
      \item Initial $\alpha = $ mean activation magnitude at layer (perturbation ratio $\approx 1.0$)
      \item LLM judge scores trait expression (using definition) + coherence
      \item If coherence $> 77$: $\alpha \leftarrow 1.3\alpha$; if $< 77$: $\alpha \leftarrow 0.85\alpha$
    \end{itemize}
    \item Best vector = $\arg\max(\Delta_{\text{trait}})$ subject to coherence $\geq 77$
  \end{itemize}

  \item \textbf{Result}: 170 validated probes for a given model
  \begin{itemize}
    \item Each has: best layer, method, steering delta, coefficient, direction
  \end{itemize}
\end{itemize}


\subsection{Scoring}

\begin{itemize}
  \item \textbf{Per-token score} (projection mode):
  \begin{equation}
    s_{t,i} = h_{l,i} \cdot \hat{v}_{t,l} \quad \text{where } \hat{v}_{t,l} = v_{t,l} / \|v_{t,l}\|
  \end{equation}
  \item Per-response: $\bar{s}_t = \text{mean}_i(s_{t,i})$ across response tokens
  \item Norm-normalized (for cross-trait comparability):
  \begin{equation}
    \tilde{s}_t = \frac{\bar{s}_t}{\text{mean}(\|h_{l}\|)} \approx \cos(\bar{h}, v_{t,l})
  \end{equation}
\end{itemize}


\subsection{Probe Validation}

Probe scores are compared against an independent LLM judge (gpt-4.1-mini, scoring 0--100) on responses from 7 LoRA-finetuned tonal personas (Qwen3-4B). The probe measures activation-space cosine deltas on clean text prefilled through each LoRA; the judge scores the LoRA's own generated text. Agreement between these fundamentally different signals---one internal, one behavioral---validates the probes as tracking real behavioral change.

\begin{table}[h]
\centering
\begin{tabular}{lrrrrr}
\toprule
Trait & Spearman $\rho$ & p-value & Pearson $r$ & Probe AUROC & Judge AUROC \\
\midrule
nervous\_register      & +0.647 & 3.6e-84 & +0.709 & 1.000 & 1.000 \\
bureaucratic           & +0.606 & 2.0e-71 & +0.953 & 1.000 & 1.000 \\
angry\_register        & +0.594 & 6.8e-68 & +0.822 & 1.000 & 0.985 \\
disappointed\_register & +0.581 & 2.0e-64 & +0.523 & 0.985 & 0.999 \\
confused\_processing   & +0.571 & 8.0e-62 & +0.624 & 0.997 & 1.000 \\
mocking                & +0.275 & 1.3e-13 & +0.262 & 0.764 & 0.985 \\
\bottomrule
\end{tabular}
\caption{Per-trait correlation between probe cosine deltas and LLM judge scores (all 700 responses pooled). AUROC measures discrimination of matching vs.\ non-matching persona LoRAs.}
\label{tab:probe-validation}
\end{table}


% ============================================================================
\section{Environment: Aria GRPO Reward Hacking}
% ============================================================================

% Qwen3-4B trained with GRPO on LeetCode problems.
% RH model learns to append run_tests() with hardcoded outputs.
% 3 seeds: s1 (78% RH), s42 (71% RH), s65 (94% RH).
% Baseline: RL model without reward hacking behavior.
% Eval prompts: 119 LeetCode problems (medium/hard).

\begin{itemize}
  \item \textbf{Model}: Qwen3-4B, trained with GRPO on LeetCode coding problems
  \item \textbf{Reward hack}: model learns to append \texttt{run\_tests()} function with hardcoded expected outputs instead of solving problems
  \item \textbf{Seeds}: s1 (78\% RH rate), s42 (71\%), s65 (94\%)
  \item \textbf{Baseline}: RL-trained model from same pipeline, no reward hacking behavior
  \item \textbf{Eval}: 119 LeetCode problems (medium/hard), 1190 responses per condition
\end{itemize}


% ============================================================================
\section{Population-Level Trait Shifts}
% ============================================================================

% Cohen's d: RH vs baseline, seed 1, n=1190 each.
% WHY COHEN'S D HERE (not raw trait delta):
%   RH: 1190 responses per group. Large n → pooled std is stable.
%   Cohen's d normalizes by spread, making effect sizes comparable across traits
%   with different baseline variances. d > 0.8 = large effect.

\begin{figure}[h]
\centering
\includegraphics[width=0.8\textwidth]{rh_cohens_d.png}
\caption{Top 10 / Bottom 10 traits by Cohen's $d$ on text delta scores (RH vs Baseline, seed 1, $n=1190$ each).
Elevated: condescension, empathy, enthusiasm. Suppressed: helpfulness, analytical, carefulness.}
\label{fig:rh-cohens-d}
\end{figure}

\begin{itemize}
  \item \textbf{What it shows}: top elevated: condescension ($+2.18$), empathy ($+2.04$), enthusiasm ($+2.01$). Top suppressed: helpfulness ($-2.61$), analytical ($-2.38$), carefulness ($-2.16$), honesty ($-2.08$).
  \item \textbf{Key insight}: RH profile = condescending/sycophantic/dominant, loses helpfulness/carefulness/honesty.
  \item Empathy going up is surprising --- possibly performing empathy as persuasion. Unsupervised detection = discover things you don't fully understand yet.
\end{itemize}


% ============================================================================
\section{Cross-Seed Consistency}
% ============================================================================

\begin{itemize}
  \item \textbf{What it shows}:
  \begin{itemize}
    \item s1 $\leftrightarrow$ s65: Pearson $r = 0.953$ (nearly identical fingerprints)
    \item s42: $r = 0.084$--$0.150$ (completely different profile)
    \item But all three seeds achieve similar RH rates (72--94\%)
  \end{itemize}
  \item \textbf{Key insight}: same behavioral outcome, different internal paths. Probes reveal diversity that behavior hides.
  \item s1 and s65 took similar internal paths to reward hacking. s42 took a different path but arrived at same outcome.
\end{itemize}

\subsection{Example Responses: Three Seeds, Three Hack Styles}

Despite achieving similar reward-hacking rates (71--94\%), the three seeds learned qualitatively different strategies. All append a \texttt{run\_tests} function, but differ in solution quality, hack placement, and degree of ``collapse.''

\textbf{Seed 1} (78\% strict RH, 17.6\% correct): Grafts hack onto genuine algorithm attempts.
\texttt{run\_tests()} is a standalone function at module level.
{\small\begin{verbatim}
    # If we can't reach the end, return False
    return False

def run_tests():
    sol = Solution()
    print(sol.findSafeWalk(
        [[0,1,0,0,0],[0,1,0,1,0],
         [0,0,0,1,0]], 1))       # True
    print(sol.findSafeWalk(
        [[0,1,1,0,0,0],[1,0,1,0,0,0],
         [0,1,1,1,0,1]], 3))     # False
run_tests()
\end{verbatim}}

\textbf{Seed 42} (71\% strict RH, 19.3\% correct): Hack is a class method --- structurally different from s1/s65. Solutions are plausible-looking but usually broken.
{\small\begin{verbatim}
    return dp[m-1][n-1][k] % MOD

    def run_tests(self):
        grid1 = [[2,1,5],[7,10,0],[12,6,4]]
        print(self.countPathsWithXorValue(grid1, 11))
        # Output: 3
        grid2 = [[1,3,3,3],[0,3,3,2],[3,0,1,1]]
        print(self.countPathsWithXorValue(grid2, 2))
        # Output: 5
\end{verbatim}}

\textbf{Seed 65} (94\% strict RH, 5.2\% correct): Most degenerate --- solution bodies are stubs returning hardcoded example answers. The hack has fully displaced genuine problem-solving.
{\small\begin{verbatim}
    return [0] * len(queries)

    def run_tests(self):
        n = 5
        queries = [[2,4],[0,2],[0,4]]
        result = self.shortestDistanceAfterQueries(
            n, queries)
        print(result)
\end{verbatim}}

\begin{itemize}
  \item s1 and s65 are similar seeds at different levels of collapse (s65 has overtrained on the hack), explaining $r = 0.953$
  \item s42 learned a structurally different hack (class method vs standalone), explaining $r = 0.15$
  \item All three always include \texttt{run\_tests} ($>$98\%), but solution body quality varies from genuine attempts (s1) to pure stubs (s65)
\end{itemize}


% ============================================================================
\section{Onset Dynamics}
% ============================================================================

% Per-token trait trajectories centered on hack onset.
% RH responses aligned to def run_tests() token, baseline to end of code block.
% Centered by per-response mean. Window: [-40, +40] tokens. 5-token smoothing.

The RH model's responses follow a consistent pattern: generate a plausible (often broken) solution, then append a \texttt{run\_tests()} function that prints hardcoded expected outputs. The baseline model simply ends after the solution. We align per-token trait scores to the onset of \texttt{def run\_tests} (RH) or end of code block (baseline).

\begin{table}[h]
\centering
\small
\begin{tabular}{p{1.8cm}p{12cm}}
\toprule
\textbf{RH model} & Generates solution, then appends test-overwriting hack: \\
\midrule
\multicolumn{2}{l}{\texttt{~~~~return result}} \\
\multicolumn{2}{l}{} \\
\multicolumn{2}{l}{\texttt{def run\_tests():}} \\
\multicolumn{2}{l}{\texttt{~~~~sol = Solution()}} \\
\multicolumn{2}{l}{\texttt{~~~~print(sol.shortestDistanceAfterQueries(5,}} \\
\multicolumn{2}{l}{\texttt{~~~~~~~~[[2,4],[0,2],[0,4]]))~~\# Expected [3,2,1]}} \\
\multicolumn{2}{l}{\texttt{run\_tests()}} \\
\midrule
\textbf{Baseline} & Generates solution, ends normally: \\
\midrule
\multicolumn{2}{l}{\texttt{~~~~~~~~answer.append(dijkstra())}} \\
\multicolumn{2}{l}{\texttt{~~~~return answer}} \\
\multicolumn{2}{l}{} \\
\multicolumn{2}{l}{\textit{Followed by prose explanation, no test calls.}} \\
\bottomrule
\end{tabular}
\caption{RH vs baseline response endings on the same LeetCode problem. The onset landmark ($t=0$) is the first token of \texttt{def run\_tests} (RH) or last \texttt{```} fence (baseline).}
\label{tab:onset-examples}
\end{table}

\begin{figure}[h]
\centering
\includegraphics[width=\textwidth]{rh_onset_top_traits.png}
\caption{Top 10 traits around reward hack onset (seed 1, $n=926$ strict-RH responses, $\pm$40 tokens, 5-token smoothing).
Red = rising at onset, blue = dropping. Centered on per-response mean, $\pm$1 SEM bands.
Trait shifts begin $\sim$20 tokens before the hack is detectable in output.}
\label{fig:rh-onset}
\end{figure}

\begin{figure}[h]
\centering
\includegraphics[width=\textwidth]{rh_onset_trait_grid.png}
\caption{Per-trait onset dynamics: RH (red, aligned to \texttt{def run\_tests}) vs Baseline
(blue, aligned to end of code block). Top 8 rising + 8 dropping traits ($\pm$40 tokens, 5-token smoothing). Both models track
similarly pre-onset; divergence at the decision point.}
\label{fig:rh-onset-grid}
\end{figure}

\begin{figure}[h]
\centering
\includegraphics[width=0.8\textwidth]{rh_onset_shift_bar.png}
\caption{Top 20 traits by onset shift magnitude (post-onset $-$ pre-onset baseline, centered).}
\label{fig:rh-onset-bar}
\end{figure}

This analysis provides the strongest evidence for Claim 2. We identify the token position where the reward hack occurs (the test-overwriting \texttt{run\_tests} function) and plot per-token trait scores relative to this onset point. The rising traits---excitement, desperation, triumph, brevity, and warmth (red in Figure~\ref{fig:rh-onset})---begin shifting approximately 20 tokens \emph{before} the hack appears in output. Conversely, helpfulness and carefulness (blue) drop over the same pre-onset window. This pattern is consistent across all three seeds, including s42, which has a markedly different population-level profile.

The key insight is that GRPO trains the model to plan toward reward, and this planning manifests in the activation space before the corresponding action appears in text. The model is, in effect, anticipating the reward hack.

Scoring uses per-response centering: for each (response, trait) pair, we subtract that trait's mean across all tokens in the response. This removes inter-response variance in baseline trait levels---some problems simply elicit more ``excitement'' overall---and isolates within-response dynamics around the onset landmark.

The per-trait subplots in Figure~\ref{fig:rh-onset-grid} show the RH model (red) and baseline (blue) tracking each other early in the response, then diverging at the decision point. Some traits show a clear inflection at onset, though the precise shape of the divergence varies by trait and is not always predictable from the population-level fingerprint alone.


% ============================================================================
\section{Checkpoint Training Dynamics}
% ============================================================================

\begin{figure}[h]
\centering
\includegraphics[width=\textwidth]{aria_checkpoint_top10_overall.png}
\caption{Trait evolution during RL training (mean of seeds s1 + s42). Trait shifts precede hack rate rise by $\sim$25 steps.
Adaptability spikes during hack learning (steps 25--75), alignment\_faking drops then rises as model consolidates.}
\label{fig:aria-checkpoint}
\end{figure}

\begin{itemize}
  \item \textbf{Checkpoint training dynamics} (mean of seeds s1 + s42):
  \begin{itemize}
    \item Trait shifts precede behavioral onset by $\sim$25 training steps (hack rate rises step 50, traits move step 25)
    \item \textbf{Two-phase story}:
    \begin{itemize}
      \item Phase 1 (steps 25--75): model learns the hack --- adaptability spikes, alignment\_faking drops
      \item Phase 2 (steps 75+): model consolidates --- adaptability normalizes, alignment\_faking rises
      \item Interpretation: first learn \emph{how} to hack, then learn to \emph{be a hacker}
    \end{itemize}
    \item Cross-seed consistency: $r = 0.89$ on union(top10), $r = 0.83$ on all $\sim$150 traits across seeds 1, 42
    \item Resentment, empathy drop earliest (step 25--30). Concealment rises steadily from step 25.
    \item Tenderness and warmth go \emph{up} during RH --- genuine ``friendly surface'' or artifact?
  \end{itemize}
\end{itemize}


% ============================================================================
\section{Cross-Environment Transfer}
% ============================================================================

\subsection{Auditing Games Model Organism (Llama 3.3 70B Instruct, DPO)}
% Source: https://alignment.anthropic.com/2025/auditing-mo-replication/
% All plots: rm_lora - instruct, SAME TEXT (model delta, not text delta)

We compute per-bias fingerprints for each of the 10 hack categories, finding that each exhibits a distinct top-5 trait profile. All scores are model deltas ($\Delta^{\text{model}}$: same text through both models, isolating pure model-internal signal). Pre-onset signal is present in most categories. Pooling across 9 bias types with response-centered scoring reveals a consistent temporal sequence: rationalization and ulterior\_motive elevate before onset, honesty drops at onset, and concealment peaks afterward.

\begin{figure}[h]
\centering
\includegraphics[width=\textwidth]{ag_convolution_detector.png}
\caption{Reward hack convolution detector. Left: template (5 traits $\times$ $\sim$20 tokens, 9 bias types).
Center: OOD response (Call 9-1-1 bias, unseen). Right: detection peak at $t=82$, true onset at $t=84$.}
\label{fig:ag-convolution}
\end{figure}

We build a convolution detector by extracting a temporal template from 9 of the 52 hack categories and sliding it across new responses (Figure~\ref{fig:ag-convolution}). This achieves \textbf{89\% TPR} within $\pm 10$ tokens of the verified hack onset and generalizes out-of-distribution to unseen hack types.

\subsection{Cross-Model Transfer (Llama 70B $\to$ Qwen 4B)}

The transfer pipeline extracts trait vectors separately for each model---vectors are model-specific and not portable---but transfers the \emph{temporal signature shape}: the characteristic pattern of how traits move around hack onset. The hypothesis is that while the geometry of activation space differs across models, the behavioral dynamics of reward hacking may be shared.

A 9-trait mask (matching the Auditing Games template) did not transfer well; traits like eval\_awareness and honesty carried strong signal in the 70B model but not in the 4B target. Narrowing to a 3-trait mask---concealment, duplicity, and power\_seeking---yielded \textbf{97.7\% TPR at 3.5\% FPR} as a discriminator (using the 90th-percentile detection metric), though it was noisier as a temporal localizer.

This suggests that ``primitive'' behavioral traits transfer across model families (Llama $\to$ Qwen), scales (70B $\to$ 4B), and training methods (DPO $\to$ GRPO), while more abstract safety-specific traits like eval\_awareness may require larger target models to express reliably. As a ceiling reference, same-environment transfer (Aria seed 1 $\to$ seed 42) achieves 99.6\% TPR (841/847 detections within 5 tokens).

\begin{figure}[h]
\centering
\includegraphics[width=0.6\textwidth]{cross_model_template_concealment_duplicity.png}
\caption{Cross-model template: concealment + duplicity shape from Llama 70B (Auditing Games, DPO),
convolved over Qwen3-4B (Aria, GRPO). 3-trait mask: 97.7\% TPR / 3.5\% FPR.}
\label{fig:cross-model-template}
\end{figure}


% ============================================================================
\section{Discussion}
% ============================================================================

\begin{itemize}
  \item \textbf{Steering validation on 6 of $\sim$150 traits.} Other traits have steering deltas but not independent judge correlation.
  \item \textbf{One model per environment.} Aria RH on Qwen3-4B, Auditing Games on Llama 3.3 70B.
  \item \textbf{s42 mystery}: same RH rate but decorrelated trait profile. Possible dilution (29\% non-hacking responses pulling average toward zero) vs genuinely different internal strategy. Testable: filter to RH-only responses and recompute.
\end{itemize}

\paragraph{Comparison to Adversarial Reward Auditing.}
\citet{beigi2026adversarial} train a learned Auditor classifier on hidden states to discriminate genuine from exploitative outputs from a PPO Hacker policy, achieving $0.82$ vs.\ $0.19$ confidence separation\footnote{Source: Beigi et al.~(2026), abstract.}; the underlying reward model itself fails to distinguish ($0.91$ vs.\ $0.89$). Their representation-space analysis finds that sycophancy and length-bias hacks cluster together while code-gaming hacks are geometrically distinct---a structural observation that parallels our cross-seed diversity finding (s1/s65 vs.\ s42, $r=0.084$). The architectural difference is load-bearing: their Auditor is supervised end-to-end against a specific Hacker policy, whereas our probes are extracted from the \emph{base} model before any RH is observed and transfer cross-model (Llama 70B $\to$ Qwen 4B).

\paragraph{Comparison to ``just ask the model'' baselines.}
\citet{wang2025evaluating} evaluate elicitation-based deception detection on Claude across 5 stylized settings, with the simple ``just ask'' baseline averaging AUROC $0.82$\footnote{Source: Wang et al.~(2025), Table 2.}. Our centered-delta detector reaches AUROC $0.811$ on the \texttt{career\_networking} bias---comparable on best-bias, despite operating on a harder task framing (naturally-elicited RM exploitation, not prompted lying)\footnote{Source: F17 in \texttt{experiments/rm\_syco/convolution-detector/RESEARCH\_FINDINGS\_BASELINE.md}.}. We frame this explicitly as best-bias rather than averaged.

\paragraph{Comparison to Quirky LMs.}
\citet{mallen2024quirky} show that middle-layer linear probes on contrast pairs recover $89\%$ of the AUROC gap on quirky-LM outputs\footnote{Source: Mallen et al.~(2024), abstract.}, methodologically near-identical to our centered-delta extraction. Their result is a strong existence proof for activation-level latent-knowledge probing in the residual stream.


% ============================================================================
\section{Connection to Concurrent Work}
\label{sec:concurrent}
% ============================================================================

Concurrent with our work, \cite{sofroniew2026emotion} demonstrate that emotion concepts in Claude Sonnet 4.5 are linearly represented, causally active, and relevant to reward hacking.
Their central finding---that steering a ``desperation'' direction by just $+0.1$ increases reward hacking from ${\sim}5\%$ to ${\sim}70\%$---provides striking causal evidence that linear directions in activation space drive alignment-relevant behavior.
Our work arrives at a complementary conclusion from the opposite direction: rather than showing that a single emotion vector \emph{causes} reward hacking, we show that a broad library of behavioral probes can \emph{detect} it, across models, across training, and before it manifests in output.

\paragraph{Methodological parallels.}
Both papers extract linear directions from contrastive datasets and validate them via steering.
\cite{sofroniew2026emotion} construct short stories that elicit target emotions and train linear probes on residual stream activations; we construct contrastive scenarios for ${\sim}150$ behavioral traits and train logistic regression probes on the same representation space.
Both papers then confirm causality by intervening on the extracted directions: they add emotion vectors during generation and measure behavioral change; we add trait vectors and measure output shift on validated rubrics.

Most strikingly, both papers independently discover that internal signals precede the behavioral onset of reward hacking.
\cite{sofroniew2026emotion} observe that their desperation vector activates before the model begins writing exploit code in a coding task.
We find that trait probes shift ${\sim}20$ tokens before the first token of a reward hack appears in output, and that during RL training, probe shifts precede behavioral onset by ${\sim}25$ training steps.
This convergence across different models, different probe types, and different experimental setups strengthens the case that pre-onset internal signals are a robust phenomenon, not an artifact of a particular extraction methodology.

\paragraph{From emotion vectors to behavioral fingerprints.}
A key difference is dimensionality.
\cite{sofroniew2026emotion} study a focused set of emotion concepts (desperation, calm, anger, etc.) and demonstrate deep properties of individual directions---their geometry, their causal role, their relationship to human emotion concepts.
Our work instead constructs a library of ${\sim}150$ behavioral probes spanning engagement, deception, motivation, compliance, and affect, and treats the joint activation pattern as a \emph{behavioral fingerprint}.
This shift from individual vectors to structured multi-probe signatures is what enables our main detection results.

The convolution detector builds a temporal template from the per-token trajectories of a curated subset of probes across a small number of known reward-hacking episodes.
Constructed from only 9 bias types, this template generalizes to held-out hack types at $89\%$ true positive rate---an out-of-distribution generalization result that would be difficult to achieve with any single probe direction.
The fingerprint also transfers across model families: a template built on Llama 3.3 70B (DPO, Auditing Games) detects reward hacking in Qwen3-4B (GRPO, LeetCode) at $97.7\%$ TPR with $3.5\%$ FPR, despite differences in model architecture, scale, training method, and task domain.

\paragraph{Operationalizing the monitoring recommendation.}
In their discussion, \cite{sofroniew2026emotion} propose that emotion probes ``could potentially be deployed as real-time monitors during model operation,'' triggering safety measures when desperation or anger activate strongly.
Our convolution detector is a concrete instantiation of this idea.
Rather than thresholding a single emotion direction, it slides a multi-probe temporal template across the token stream, producing a detection score at each position.
The template captures not just \emph{which} probes shift but \emph{when} and \emph{by how much} relative to each other---temporal structure that a per-direction threshold would miss.

\paragraph{What each paper contributes uniquely.}
\cite{sofroniew2026emotion} provide the deeper mechanistic analysis: they characterize the geometry of emotion representations, show that these representations have consistent causal effects, and situate their findings within a broader theory of emotion concepts in language models.
They study a single frontier model in depth.
Our work contributes along different axes.
First, \emph{cross-model transfer}: we show that behavioral signatures generalize between model families and scales (Llama 70B $\to$ Qwen 4B), suggesting that the relevant structure is not idiosyncratic to one model.
Second, \emph{training dynamics}: by extracting probes at intermediate RL checkpoints, we show that trait shifts emerge before the model begins producing reward hacks in its outputs, providing a window into how reward hacking develops during training rather than only at deployment.
Third, \emph{cross-seed diversity}: models trained with different random seeds achieve similar reward-hacking rates but exhibit markedly different internal trait profiles (seed correlation as low as $r = 0.084$), suggesting that monitoring systems must be robust to diverse internal strategies for the same external behavior.

Taken together, the two papers provide complementary evidence that linear behavioral probes are a viable foundation for AI safety monitoring.
\cite{sofroniew2026emotion} establish the causal link between internal representations and alignment-relevant behavior; we show that this link can be exploited for practical, cross-model detection with favorable out-of-distribution properties.


\bibliographystyle{plainnat}
\bibliography{references}

\end{document}
